\begin{pastexam}[label={px:ust-114-ee-ds-27}]{臺聯大電機類 114 學年度 問答題第 3 題}
（問答；10 分）

Based on the following code, please answer the questions below:
\begin{enumerate}[label=(\Alph*)]
\item Why can the hasCycle function successfully detect the presence of a cycle in the linked list
when link node4.next to node2? Please explain in detail how the slow and fast pointers work and at
which node they meet. (5 point)
\item If the assignment link node4.next to node2 is removed and the program is executed, what will be
the output of the hasCycle function? Please explain the reason for this situation and how the
algorithm handles it. (5 point)
\end{enumerate}

\begin{minted}[fontsize=\small]{text}
Class ListNode
    Attributes: value, next

    Constructor(value):
        value = value
        next = null

Function hasCycle(head):
    If head is null or head.next is null:
        Return false

    slowPointer = head
    fastPointer = head

    While fastPointer is not null and fastPointer.next is not null:
        slowPointer = slowPointer.next
        fastPointer = fastPointer.next.next

        If slowPointer == fastPointer:
            Return true

    Return false

Main:
    Create node1 with value 1
    Create node2 with value 2
    Create node3 with value 3
    Create node4 with value 4

    Link node1.next to node2
    Link node2.next to node3
    Link node3.next to node4
    Link node4.next to node2

    result = hasCycle(node1)
    Print "Does the linked list have a cycle? " + result
\end{minted}
\end{pastexam}

\begin{pastsoln}
串列是 node1 $\to$ node2 $\to$ node3 $\to$ node4 $\to$ node2（回到 node2），環前面有 $\mu = 1$ 個
節點（node1），環長 $\lambda = 3$（node2、node3、node4）。

\textbf{(A)} 快慢指標（Floyd 判環法，資結書第~5.6~節）：slow 每輪走一步，fast 每輪走兩步。
\begin{itemize}
\item 沒有環時 fast 一定先碰到 null，迴圈結束回傳 false；只要有環，fast 進環後就永遠在環裡繞，不會碰到
null。
\item slow 也進入環之後，把它看成 fast 在環上「追」slow。 每一輪 fast 比 slow 多走一步，兩者沿環方向的
距離減少 $1$。 距離一開始在 $0$ 到 $\lambda - 1$ 之間，所以最多 $\lambda - 1$ 輪就減到 $0$，兩個指標
指到同一個節點，回傳 true。 不會「跳過」對方，因為距離每次只少 $1$。
\end{itemize}
追蹤這個例子：
\begin{center}
\small
\begin{tabular}{ccccc}
\toprule
輪 & slow & fast & 說明 & 相等？ \\
\midrule
開始 & node1 & node1 & & （迴圈前不比） \\
1 & node2 & node3 & fast：node1 $\to$ node2 $\to$ node3 & 否 \\
2 & node3 & node2 & fast：node3 $\to$ node4 $\to$ node2 & 否 \\
3 & node4 & node4 & fast：node2 $\to$ node3 $\to$ node4 & 是 \\
\bottomrule
\end{tabular}
\end{center}
第 3 輪兩者在 \textbf{node4} 相遇，回傳 true，印出
``Does the linked list have a cycle? true''。 輪數 $3$ 符合資結書定理~5.4：$\lambda = 3$ 的倍數中不小於
$\max(\mu, 1) = 1$ 的最小者。

\textbf{(B)} 拿掉 node4.next = node2 之後，node4.next 是 null，串列 node1 $\to$ node2 $\to$ node3
$\to$ node4 $\to$ null 沒有環。
\begin{center}
\small
\begin{tabular}{cccl}
\toprule
輪 & slow & fast & 迴圈條件 \\
\midrule
開始 & node1 & node1 & fast、fast.next 都不是 null，進入 \\
1 & node2 & node3 & fast.next = node4，繼續 \\
2 & node3 & null & fast 是 null，離開 \\
\bottomrule
\end{tabular}
\end{center}
第 2 輪 fast 從 node3 走兩步：node4，再到 node4.next = null。 迴圈條件
「fastPointer is not null」不成立，結束，回傳 false，印出 ``Does the linked list have a cycle? false''。
這也說明迴圈條件為什麼要同時檢查 fastPointer 和 fastPointer.next：fast 一次跳兩步，先確定
fastPointer.next 不是 null，才能安全地取 fastPointer.next.next，否則會對 null 取欄位而出錯。 開頭的
「head 是 null 或 head.next 是 null 就回傳 false」處理空串列和只有一個節點（沒有自己指向自己）的情形。
整個函式時間 $O(n)$、額外空間 $O(1)$。
\end{pastsoln}
